% ██████████████████████████████████████████████████████████████████
%  CHAPTER 5 — IMPLEMENTATION
% ██████████████████████████████████████████████████████████████████
\chapter{Implementation}
\label{ch:implementation}

\section{Technology Stack}
\label{sec:impl_stack}

\begin{itemize}
  \item \textbf{Backend framework:} Django~4.2 (Python~3.12).
  \item \textbf{Database:} PostgreSQL with database-per-tenant isolation.
  \item \textbf{ML stack:} scikit-learn (\texttt{GradientBoostingClassifier}), SHAP (\texttt{TreeExplainer}), LIME (\texttt{LimeTabularExplainer}).
  \item \textbf{Plotting:} matplotlib (\texttt{Agg} backend), seaborn.
  \item \textbf{Reproducibility:} fixed \texttt{numpy} / scikit-learn random seeds (\texttt{SEED~=~42}); pinned dependency manifests \texttt{requirements.txt} and \texttt{requirements-ml.txt}.
\end{itemize}

\section{Sparsity Experiment Driver}
\label{sec:impl_experiment}

The experimental pipeline is encapsulated in a single standalone script, \texttt{thesis/scripts/sparsity\_simulator.py}, which (i) constructs the Late-Payment dataset from the raw Olist CSVs, (ii) produces the five stratified tenant slices, (iii) trains and evaluates one classifier per slice, (iv) computes SHAP and LIME explanations, (v) derives the stability and fidelity metrics defined in Chapter~\ref{ch:methodology}, and (vi) serialises all artefacts under \texttt{thesis/results/sparsity\_experiment/}. A companion plotting script, \texttt{plot\_sparsity\_summary.py}, produces the figures consumed in Chapter~\ref{ch:results}.

\section{Tier-Logic Gate}
\label{sec:impl_tier_gate}

The architectural contribution is implemented additively within the existing \texttt{LatePaymentPredictor} service class. The new \texttt{predict\_or\_heuristic} method is shown in Listing~\ref{lst:tier_gate}.

\begin{lstlisting}[language=Python, caption={Maturity-Aware AI Routing tier gate.}, label={lst:tier_gate}]
COLD_START_INVOICE_THRESHOLD = 1000

def predict_or_heuristic(self, invoice: Invoice) -> dict:
    """Tier 1 (cold-start): return platform-wide global heuristic.
       Tier 2 (mature):     return local model probability."""
    n = self._invoice_count()
    if n < self.COLD_START_INVOICE_THRESHOLD:
        h = self._load_global_heuristic()
        return {
            "mode": "global_heuristic",
            "tenant_invoice_count": n,
            "threshold": self.COLD_START_INVOICE_THRESHOLD,
            "reason": (
                "Local model is below the explanation-reliability "
                "threshold (Spearman rho < 0.7 vs reference). "
                "Falling back to platform-wide Global SHAP heuristic."
            ),
            "heuristic": h.get("heuristic", ""),
            "top_features": h.get("top_features", []),
            "invoice_id": invoice.pk,
        }
    return {
        "mode": "local_prediction",
        "tenant_invoice_count": n,
        "threshold": self.COLD_START_INVOICE_THRESHOLD,
        "invoice_id": invoice.pk,
        "late_probability": self.predict(invoice),
    }
\end{lstlisting}

\section{Global Heuristic Artefact}
\label{sec:impl_heuristic}

The cold-start heuristic is persisted as the JSON document \texttt{thesis/results/sparsity\_experiment/cold\_start\_heuristic.json}, generated by the experiment driver from Tenant~A's global SHAP attributions. The artefact is loaded once per process via the memoised classmethod \texttt{LatePaymentPredictor.\_load\_global\_heuristic}.

\section{Backward Compatibility}
\label{sec:impl_compat}

The existing \texttt{predict()} method is left unchanged; all upstream callers (the open-invoice batch scorer, the REST endpoint, and any scheduled Celery tasks) continue to operate without modification. Consumers wishing to benefit from the tier gate opt in by invoking \texttt{predict\_or\_heuristic()} instead.
